\pagebreak
\section{ECCO Data Filenames and Supporting Conventions}


ECCO Version 4 Release 6 (V4r6) follows a structured filename convention to organize its extensive datasets. The filenames encode essential information about the dataset's content, grid type, temporal resolution, and version. Below is an overview of the naming convention.

\subsection{General filename Structure}
\par \vspace{0.25cm}

\par For ECCO dataset product files:

\par \vspace{0.5cm}
\begin{center}
\small{\textbf{\fontfamily{lmtt}\selectfont{[ShortName]\_[TemporalResolution]\_[Indicative Time]\_ECCO\_[Version]\_[GridType].<File Type>}}}
\end{center}
\par \vspace{0.5cm}
\par For ECCO data geometry and grid files:

\par \vspace{0.5cm}
\begin{center}
    \small{\textbf{\fontfamily{lmtt}\selectfont{[ShortName]\_ECCO\_[Version]\_[GridType].<File Type>}}}
\end{center}
% \texttt
\par \vspace{0.25cm}

\begin{center}
\begin{tabular}{m{0.18\textwidth} m{0.75\textwidth}}
    \multicolumn{1}{c}{\textbf{Keypoins}} & \multicolumn{1}{c}{\textbf{Description}} \\ \hline
    ShortName: & Describes the dataset variable (e.g., SSH for sea surface height, TEMP\_SALINITY for temperature and salinity). \\ \hline 
    TemporalResolution: & Indicates the time averaging or snapshot type (e.g., DAILY, MONTHLY, SNAP). \\ \hline
    GridType : & Specifies the grid used (e.g., \textbf{native\_llc90} for the native ECCO grid or \textbf{laton\_0p50deg} for interpolated 0.5 \textdegree lat-lon grid). It can also be a specific information such as \textbf{1D} to indication the one-dimentional dataset files over the full period of ECCO data availability.\\ \hline
    Version: &  Identifies the ECCO release version (e.g., V4r6). \\ \hline
    Indicative Time: & Encodes the time period covered by the file (e.g., monthly files use YYYY-MM, daily files use YYYY-MM-DD) \\ \hline
    File Type & netCDF or Zarr data format. \\ \hline
\end{tabular}
\end{center}

\par \vspace{0.25cm}
\subsection{Examples}
\par Daily atmosphere surface temperature, humidity, winds, and pressure on the lat-lon-cap 90 (llc90)
native model grid (first) and on the interpolated regular 0.5-degree lat-lon grid (second) from ECCO V4r6 on 2017-12-29.
\begin{itemize}
    \item ATM\_SURFACE\_TEMP\_HUM\_WIND\_PRES\_day\_mean\_2017-12-29\_ECCO\_V4r6\_native\_llc0090.nc
    \item ATM\_SURFACE\_TEMP\_HUM\_WIND\_PRES\_day\_mean\_2017-12-29\_ECCO\_V4r6\_latlon\_0p50deg.nc
\end{itemize}

\par Dynamic sea surface height interpolated on the lat-lon regular 0.5-degree model grid from ECCO V4r6 on 2017-12-29.

\begin{itemize}
    \item SEA\_SURFACE\_HEIGHT\_day\_mean\_2017-12-29\_ECCO\_V4r6\_native\_llc0090.nc
\end{itemize}

\par One-dimentional field of global mean atmospheric surface pressure over the ocean and sea-ice fromthe ECCO V4r6.
\begin{itemize}
    \item GLOBAL\_MEAN\_ATM\_SURFACE\_PRES\_snap\_ECCO\_V4r6\_1D.nc
\end{itemize}

\par Geometric parameters for the regular 0.5-degree lat-lon grid from ECCO V4r6.
\begin{itemize}
    \item GRID\_GEOMETRY\_ECCO\_V4r6\_lat-lon\_0p50deg.nc
\end{itemize}
